% Q003 · 一维无限深势阱的能级跃迁与光子频率
一维无限深势阱的定态能级为
\[
E_n=\frac{n^2\pi^2\hbar^2}{2mL^2}.
\]
由 $n=2$ 向 $n=1$ 跃迁时，能量差为
\[
\Delta E=E_2-E_1
=\frac{(4-1)\pi^2\hbar^2}{2mL^2}
=\frac{3\pi^2\hbar^2}{2mL^2}>0.
\]
假设跃迁的能量全部由一个光子带走，且忽略反冲，则光子满足 $h\nu=\Delta E$。由于 $h=2\pi\hbar$，故
\[
\boxed{\nu=\frac{\Delta E}{h}
=\frac{3\pi^2\hbar^2}{2mL^2h}
=\frac{3\pi\hbar}{4mL^2}}.
\]
这表明固定粒子质量时，跃迁频率与阱宽的平方成反比。右式量纲为 $\mathrm{s^{-1}}$。
